\subsection{Component-Adaptive Tversky Term}

Lesion segmentation in brain MRI is characterized by severe class imbalance
and large variability in lesion size. Standard overlap-based losses, such as
the Dice loss or the Tversky loss~\cite{salehi_tversky_2017}, operate at the
voxel level and therefore tend to be dominated by large lesions. As a result,
small lesions may contribute negligibly to the optimization objective,
leading to reduced detection sensitivity for clinically relevant small
structures.

To mitigate this issue, we introduce a \emph{Component-Adaptive Tversky (CAT)}
term that balances lesion contributions by incorporating connected-component
information from the ground-truth segmentation.

\paragraph{Notation.}
Let $\Omega \subset \mathbb{R}^3$ denote the voxel domain of a 3D MRI volume.
For each voxel $i \in \Omega$:

\begin{itemize}
    \item $g_i \in \{0,1\}$ denotes the ground-truth label,
    \item $p_i \in [0,1]$ denotes the predicted probability of the foreground class.
\end{itemize}

The set of foreground voxels is defined as

\[
G = \{ i \in \Omega \mid g_i = 1 \}.
\]

\paragraph{Connected-component decomposition.}
The binary foreground mask is decomposed into $K$ disjoint connected components

\[
G = \bigcup_{k=1}^{K} C_k, 
\qquad
C_k \cap C_j = \varnothing \;\; \text{for } k \neq j,
\]

where each $C_k$ represents an individual lesion instance.
The size of lesion $k$ is

\[
|C_k| = \sum_{i \in C_k} 1 .
\]

\paragraph{Component-adaptive weighting.}
To balance lesion contributions during optimization,
we introduce a voxel-wise weight $w_i$ defined as

\[
w_i =
\begin{cases}
(|C_k| + \epsilon)^{-\gamma}, & \text{if } i \in C_k, \\
w_{\mathrm{bg}}, & \text{if } g_i = 0 ,
\end{cases}
\]

where $\gamma \ge 0$ controls the strength of size adaptation,
$\epsilon > 0$ stabilizes extremely small components,
and $w_{\mathrm{bg}}$ denotes the background weight.

This formulation assigns larger weights to voxels belonging to smaller
lesions and smaller weights to voxels belonging to larger lesions,
thereby reducing the dominance of large lesions in the optimization
objective.

\paragraph{Weighted Tversky index.}
Building upon the Tversky similarity index~\cite{salehi_tversky_2017},
we define weighted true positives, false positives, and false negatives as

\[
\mathrm{TP}_w = \sum_{i \in \Omega} w_i \, p_i \, g_i,
\]

\[
\mathrm{FP}_w = \sum_{i \in \Omega} w_i \, p_i \, (1 - g_i),
\]

\[
\mathrm{FN}_w = \sum_{i \in \Omega} w_i \, (1 - p_i) \, g_i.
\]

The component-adaptive Tversky index is then

\[
\mathrm{I}_{\mathrm{CAT}} =
\frac{\mathrm{TP}_w + \delta}
{\mathrm{TP}_w + \alpha \mathrm{FP}_w + \beta \mathrm{FN}_w + \delta},
\]

where $\alpha$ and $\beta$ control the penalties for false positives and
false negatives, respectively, and $\delta > 0$ is a smoothing constant.

\paragraph{CAT term.}
The final term is defined as

\[
\mathcal{L}_{\mathrm{CAT}} = 1 - \mathrm{I}_{\mathrm{CAT}} .
\]

\paragraph{Lesion-level interpretation.}
For $\gamma = 1$, the total contribution of lesion $k$ is approximately

\[
\sum_{i \in C_k} w_i
\approx
|C_k| (|C_k| + \epsilon)^{-1}
\approx 1,
\]

indicating that each lesion contributes approximately equally to the
optimization objective regardless of its volume. This shifts the learning
process from voxel-dominated optimization toward lesion-balanced
optimization, which is particularly beneficial for improving the
segmentation of small lesions.

\paragraph{Relationship to existing losses.}
The proposed formulation generalizes several existing overlap-based losses:

\begin{itemize}
\item If $\gamma = 0$, the CAT term reduces to the standard Tversky loss.
\item If $\alpha = \beta = 0.5$, the formulation corresponds to a component-weighted Dice loss.
\item If $w_i = 1$ for all voxels, the formulation becomes the standard voxel-wise Tversky loss.
\end{itemize}

\subsection{Lesion-Level Multiple Instance Learning Term}

While overlap-based losses encourage accurate voxel-wise segmentation,
they do not explicitly enforce the detection of individual lesions.
In highly imbalanced medical segmentation tasks, a model may achieve
reasonable overlap metrics while completely missing small lesions.
To address this limitation, we introduce an auxiliary
\emph{lesion-level detection term} based on the Multiple Instance
Learning (MIL) paradigm.

\paragraph{Lesion instances.}

Let the foreground set $G$ be decomposed into $K$ connected components

\[
G = \bigcup_{k=1}^{K} C_k,
\]

where each component $C_k$ represents an individual lesion instance.

\paragraph{Instance detection score.}

Let $p_i \in [0,1]$ denote the predicted foreground probability for voxel
$i$. For each lesion component $C_k$, we define the detection score as

\[
s_k = \max_{i \in C_k} p_i .
\]

\paragraph{MIL term.}

\[
\mathcal{L}_{\mathrm{MIL}}
=
\frac{1}{K}
\sum_{k=1}^{K}
-\log \left( s_k + \epsilon \right),
\]

where $\epsilon$ is a small constant ensuring numerical stability.
For samples with no foreground lesion components (i.e., $K = 0$),
$\mathcal{L}_{\mathrm{MIL}}$ is set to $0$.

This objective encourages at least one voxel inside each lesion to have
high predicted probability, thereby improving lesion-level recall,
particularly for small lesions.

\subsection{CATMIL Objective Function}

The base loss is defined as

\[
\mathcal{L}_{\mathrm{base}}
=
\mathcal{L}_{\mathrm{Dice}}
+
\mathcal{L}_{\mathrm{CE}}.
\]

The proposed objective function, referred to as CATMIL, integrates the
component-adaptive term and the lesion-level term into a unified training
objective.

\[
\mathcal{L}_{\mathrm{CATMIL}}
=
\mathcal{L}_{\mathrm{base}}
+
\lambda_{\mathrm{CAT}}
\mathcal{L}_{\mathrm{CAT}}
+
\lambda_{\mathrm{MIL}}
\mathcal{L}_{\mathrm{MIL}},
\]

where $\mathcal{L}_{\mathrm{CAT}}$ is the component-adaptive term and
$\mathcal{L}_{\mathrm{MIL}}$ is the lesion-level detection term.

\paragraph{Training strategy.}

\[
\lambda_{\mathrm{CAT}}(t)
=
\lambda_{\mathrm{CAT}}^{\mathrm{final}}
\cdot
\min\left(\frac{t}{T},1\right),
\]

where $t$ denotes the current training epoch and $T$ is the number of
warm-up epochs. The MIL weight $\lambda_{\mathrm{MIL}}$ is kept constant
during training.